\section{Proof-of-Concept and Pilot Evaluation}
\label{sec:evaluation}

\subsection{Evaluation Design}

We conducted a proof-of-concept comparative user study contrasting a scripted baseline with the adaptive system. The study design was a within-subjects pilot comparison where participants experienced both conditions in counterbalanced order.

\subsection{Participants and Protocol}

Three participants (P1, P2, P3) aged 28–35 participated in the pilot study. Each participant completed two stretching sessions, each lasting approximately 8–10 minutes: (1) a scripted baseline condition where the robot followed a fixed exercise sequence without perception-based adaptation, and (2) the adaptive condition where the robot adjusted guidance based on user state and environment.

A two-minute break separated the two conditions. Participants were informed of the study's purpose and signed informed consent forms. The study was conducted in a controlled laboratory environment with standard objects present (coffee mug, water bottle, banana, chair).

\subsection{Quantitative Metrics}

After each condition, participants completed a custom questionnaire rating the robot's performance on six dimensions:

\begin{enumerate}
\item \textbf{Clarity:} How clear were the robot's instructions? (1–5 Likert scale)
\item \textbf{Comfort:} How comfortable were you during the routine? (1–5)
\item \textbf{Adaptability:} How well did the robot adapt to your state and needs? (1–5)
\item \textbf{Trust:} How much did you trust the robot's decisions? (1–5)
\item \textbf{Naturalness:} How natural and human-like was the robot's behavior? (1–5)
\item \textbf{Object Relevance:} How relevant were the robot's object suggestions? (1–5, N/A for scripted condition)
\end{enumerate}

Additionally, we recorded:
\begin{itemize}
\item Task completion time (total session duration).
\item Number of user interventions (corrections or clarifications requested).
\item Number of LLM reasoning cycles per session.
\item Latency of adaptive responses (time from user input to robot output).
\end{itemize}

\subsection{Qualitative Feedback}

After completing both conditions, participants participated in semi-structured interviews (15–20 minutes) discussing:
\begin{itemize}
\item Which version they preferred and why.
\item Perceived strengths and weaknesses of each approach.
\item Suggestions for improvement.
\item Overall sense of engagement and repeated-use intentions.
\end{itemize}

All interviews were recorded, transcribed, and analyzed using thematic coding by two independent raters. Inter-rater agreement (Cohen's kappa) was monitored to ensure consistency.

\subsection{Planned Analysis for a Full-Scale Study}

To move beyond pilot-level evidence, the next evaluation phase will use a pre-registered analysis plan. The primary endpoint will be perceived adaptability, with secondary endpoints including trust, naturalness, completion time, and intervention count. For subjective outcomes, we will use mixed-effects models with participant-level random intercepts to account for repeated measures. For non-normal outcomes (e.g., count-based interventions), we will use generalized mixed-effects models.

We will also report effect sizes and confidence intervals in addition to p-values, and include multiple-comparison correction for secondary outcomes. This design will allow us to test whether the trends observed in this pilot remain stable at larger scale while preserving ecological validity.
